% =====================================================================
% POPW Multi-Task Assembly Recognition --- Benchmark Comparison Template
% =====================================================================
% Compile: pdflatex popw_benchmark.tex && bibtex popw_benchmark && pdflatex popw_benchmark.tex && pdflatex popw_benchmark.tex
% Or with biber: latexmk -pdf popw_benchmark.tex
% =====================================================================

\documentclass[11pt,a4paper]{article}

% ---------- Geometry & spacing ----------
\usepackage[a4paper,margin=1in]{geometry}
\usepackage[final,protrusion=true,expansion=false]{microtype}
\usepackage[T1]{fontenc}
\usepackage[utf8]{inputenc}

% ---------- Tables ----------
\usepackage{booktabs}
\usepackage{multirow}
\usepackage{array}
\usepackage{tabularx}
\usepackage{makecell}
\usepackage{siunitx}
\sisetup{
  detect-weight        = true,
  detect-family        = true,
  group-separator      = {,},
  group-minimum-digits = 4,
  table-format         = 2.2,
  table-number-alignment = center,
}

% ---------- Math, graphics, links ----------
\usepackage{amsmath,amssymb}
\usepackage{graphicx}
\usepackage[hidelinks]{hyperref}
\usepackage{xcolor}
\usepackage{url}

% ---------- Bibliography ----------
\usepackage[numbers,sort&compress]{natbib}

% ---------- Float control ----------
\usepackage{placeins}        % \FloatBarrier
\usepackage{float}           % H placement specifier

% ---------- Convenience ----------
\newcommand{\popw}{POPW}
\newcommand{\ikea}{IKEA-ASM}
\newcommand{\industreal}{IndustReal}
\newcommand{\map}{mAP}
\newcommand{\todo}{\textit{---}}              % placeholder for unfilled results
\newcommand{\best}[1]{\textbf{#1}}            % bold = best in column
\newcommand{\runnerup}[1]{\underline{#1}}     % underline = second best
\newcommand{\popwres}{\todo}                  % shortcut for empty POPW cell

% Symbols for table footnotes
\newcommand{\sym}[1]{\textsuperscript{\textit{#1}}}

\title{\popw{}: A Unified Multi-Task Architecture for\\
       Egocentric Assembly Understanding\\
       \textit{Benchmark Comparison Tables}}
\author{(Author names omitted for review)}
\date{}

\begin{document}
\maketitle

\begin{abstract}
We present \popw{}, a single unified architecture for multi-task assembly understanding from egocentric video.
A shared ConvNeXt-Tiny backbone with a Feature Pyramid Network feeds five task heads:
RetinaNet-style object/state detection, 2D pose regression, 9-DoF head pose,
classifier-attention activity recognition, and a causal-Transformer procedure step recognition (PSR) head.
The same architecture is evaluated on \ikea{}~\cite{benshabat2021ikea} and \industreal{}~\cite{schoonbeek2024industreal},
adapted only by output dimensionality.
This document collects the benchmark tables against which we compare; all \popw{} cells are
intentionally left blank during the experimental phase.
\end{abstract}

% =====================================================================
\section{Introduction}
\label{sec:intro}

Assembly understanding is fragmented across single-task models: dedicated detectors
for object/state detection (e.g.\ YOLOv8m on \industreal{}~\cite{schoonbeek2024industreal}),
dedicated transformers for action recognition (e.g.\ MViTv2~\cite{schoonbeek2024industreal} pretrained on Kinetics),
and dedicated rule-based or temporal-encoder pipelines for procedure step recognition (PSR)
such as the \industreal{} B3 baseline and STORM-PSR~\cite{schoonbeek2025storm}.
\popw{} replaces this stack with a single forward pass, which is a meaningful systems contribution
when streaming inference latency matters.

The remainder of this document is the \emph{benchmark comparison harness}.
Tables~\ref{tab:ikea-headline}--\ref{tab:efficiency} list every external baseline
\popw{} is compared against, along with the protocol and source paper.
The \popw{} columns are deliberately blank until experiments complete;
once we have results, only those cells need to be filled.

% =====================================================================
\section{Datasets and protocols}
\label{sec:datasets}

\paragraph{\ikea{}~\cite{benshabat2021ikea}.}
371 videos, 33 atomic action classes, 17K action instances, $\sim$3M frames.
Three RGB views (front / top / side) and depth at $640{\times}480$.
Tasks: object segmentation (7 classes), 2D body pose (17 keypoints), action recognition,
phase classification, temporal localization.

\paragraph{\industreal{}~\cite{schoonbeek2024industreal}.}
Egocentric RGB at $1280{\times}720$, single camera, real and synthetic recordings of
toy-assembly tasks. 74 atomic action classes, 24 ASD (Assembly State Detection) classes,
11-component PSR labels, and 9-DoF head pose ground truth.

\paragraph{Notation.}
The \popw{} architecture is identical between datasets; only head dimensions and the
applicable task set differ. Resolution, class count, and presence of PSR / phase /
temporal-localization labels are summarised in Table~\ref{tab:dataset-config}.

\begin{table}[!htbp]
\centering
\small
\caption{Dataset configuration. \popw{} uses the same backbone, FPN, and head architectures
on both datasets; only the output dimensionality and the active task set differ.}
\label{tab:dataset-config}
\begin{tabular}{l c c}
\toprule
& \ikea{} & \industreal{} \\
\midrule
Resolution                      & $640{\times}480$ & $1280{\times}720$ \\
Cameras                         & 3 RGB views      & 1 RGB (egocentric) \\
Detection classes               & 7 objects        & 24 ASD states \\
Pose                            & body, 17 kpts    & head, 9-DoF \\
Action classes                  & 33               & 74 \\
PSR                             & ---              & 11 components \\
Phase classification            & 12-phase         & --- \\
Temporal localization           & yes              & --- \\
\bottomrule
\end{tabular}
\end{table}

\FloatBarrier
% =====================================================================
\section{Headline benchmark tables}
\label{sec:headline}

Tables~\ref{tab:ikea-headline} and~\ref{tab:industreal-headline} list every
metric for which a published baseline exists on \ikea{} and \industreal{}
respectively. Each row records the protocol, the strongest baseline we are
aware of, and the \popw{} entry under the same protocol. \popw{} entries are
deliberately blank during the experimental phase; once results converge they
will be filled in with mean$\pm$std across three random seeds.

% --------------------------- IKEA ASM ---------------------------------

\begin{table}[!htbp]
\centering
\small
\caption{\textbf{\ikea{} headline benchmarks.} Each row is a published baseline for a specific task.
\popw{} (last group) is evaluated under the same protocol as the corresponding row.
Empty cells will be filled once experiments converge; results will be reported as
mean$\pm$std across three random seeds.}
\label{tab:ikea-headline}
\renewcommand{\arraystretch}{1.15}
\begin{tabular}{l l l c c}
\toprule
\textbf{Task} & \textbf{Method} & \textbf{Metric} & \textbf{Score} & \textbf{Reference} \\
\midrule
\multicolumn{5}{l}{\textit{Object segmentation}} \\
Front view & ResNeXt-101-FPN  & AP@0.5         & 85.3 & \cite{benshabat2021ikea} \\
Front view & ResNeXt-101-FPN  & AP (COCO)      & 65.9 & \cite{benshabat2021ikea} \\
Front view & Mask R-CNN       & AP@0.5         & 78.9 & \cite{benshabat2021ikea} \\
Front view & \popw{} (ours)   & AP@0.5         & \popwres & --- \\
Front view & \popw{} (ours)   & AP (COCO)      & \popwres & --- \\
\midrule
\multicolumn{5}{l}{\textit{2D body pose}} \\
Front view & MaskRCNN-ft      & PCK@10\,px     & 64.3 & \cite{benshabat2021ikea} \\
Front view & MaskRCNN-ft      & PCK@0.2        & 88.0 & \cite{benshabat2021ikea} \\
Front view & \popw{} (ours)   & PCK@10\,px     & \popwres & --- \\
Front view & \popw{} (ours)   & PCK@0.2        & \popwres & --- \\
\midrule
\multicolumn{5}{l}{\textit{Activity recognition}} \\
Front view (RGB)        & P3D                  & Top-1   & 60.40 & \cite{benshabat2021ikea} \\
Front view (RGB+pose)   & I3D RGB+pose         & Top-1   & 64.15 & \cite{benshabat2021ikea} \\
All views (RGB)         & I3D                  & Top-1   & 47.00 & \cite{benshabat2021ikea} \\
All views, most-relevant obj. & PC3D~\cite{aganian2023ikea}     & Top-1 & 80.20 & \cite{aganian2023ikea} \\
All views, all obj.     & 2D-CNN~\cite{aganian2023ikea}        & Top-1 & 79.70 & \cite{aganian2023ikea} \\
Online (cross-subject)  & MiniROAD             & mcAP    & 80.84 & \cite{xie2025ptma} \\
Online (cross-subject)  & PTMA                 & mcAP    & 86.99 & \cite{xie2025ptma} \\
Online (cross-subj.+view) & PTMA               & mcAP    & 84.47 & \cite{xie2025ptma} \\
Front view (RGB+pose)   & \popw{} (ours)       & Top-1   & \popwres & --- \\
All views               & \popw{} (ours)       & Top-1   & \popwres & --- \\
Online                  & \popw{} (ours)       & mcAP    & \popwres & --- \\
\midrule
\multicolumn{5}{l}{\textit{Temporal localization}} \\
All views & I3D                       & \map{}@0.5 & 20.00 & \cite{benshabat2021ikea} \\
RGB-only  & ActionFormer              & \map{}@0.5 & 21.49 & \cite{gatedsrm2026} \\
RGB+pose  & Gated SRM                 & \map{}@0.5 & 21.77 & \cite{gatedsrm2026} \\
          & \popw{} (ours)            & \map{}@0.5 & \popwres & --- \\
\midrule
\multicolumn{5}{l}{\textit{Phase classification \& temporal order}} \\
Self-supervised & STEPs               & Acc@1.0          & 37.02 & \cite{shah2023steps} \\
Self-supervised & STEPs               & Kendall's $\tau$ &  0.91 & \cite{shah2023steps} \\
                & \popw{} (ours)      & Acc@1.0          & \popwres & --- \\
                & \popw{} (ours)      & Kendall's $\tau$ & \popwres & --- \\
\bottomrule
\end{tabular}
\end{table}

% --------------------------- IndustReal -------------------------------
\begin{table}[!htbp]
\centering
\small
\caption{\textbf{\industreal{} headline benchmarks.} Each baseline is the strongest competitor
reported in the cited paper. \industreal{} is single-view egocentric, so cross-view protocols
do not apply. Head-pose has no published supervised baseline; \popw{} establishes one.}
\label{tab:industreal-headline}
\renewcommand{\arraystretch}{1.15}
\begin{tabular}{l l l c c}
\toprule
\textbf{Task} & \textbf{Method} & \textbf{Metric} & \textbf{Score} & \textbf{Reference} \\
\midrule
\multicolumn{5}{l}{\textit{Assembly State Detection (ASD)}} \\
& YOLOv8m (COCO+synth+real) & \map{}@0.5     & 83.80 & \cite{schoonbeek2024industreal} \\
& \popw{} (ours)            & \map{}@0.5     & \popwres & --- \\
& \popw{} (ours)            & \map{}@[0.5{:}0.95] & \popwres & --- \\
\midrule
\multicolumn{5}{l}{\textit{Activity recognition}} \\
& MViTv2 (Kinetics-400 pretrained) & Top-1   & 66.45 & \cite{schoonbeek2024industreal} \\
& MViTv2 (Kinetics-400 pretrained) & Top-5   & 88.43 & \cite{schoonbeek2024industreal} \\
& \popw{} (ours)            & Top-1          & \popwres & --- \\
& \popw{} (ours)            & Top-5          & \popwres & --- \\
\midrule
\multicolumn{5}{l}{\textit{Procedure Step Recognition (PSR)}} \\
& B3 rule-based             & F1             & 0.883 & \cite{schoonbeek2024industreal} \\
& B3 rule-based             & Precision      & 0.885 & \cite{schoonbeek2024industreal} \\
& B3 rule-based             & Recall         & 0.880 & \cite{schoonbeek2024industreal} \\
& STORM-PSR                 & F1\sym{a}      & 0.901 & \cite{schoonbeek2025storm} \\
& STORM-PSR                 & POS\sym{a}     & 0.812 & \cite{schoonbeek2025storm} \\
& \popw{} (ours)            & F1             & \popwres & --- \\
& \popw{} (ours)            & POS            & \popwres & --- \\
\midrule
\multicolumn{5}{l}{\textit{Assembly state recognition / error verification}} \\
& SupCon+ISIL (ResNet-34)   & F1@1           & \texttt{$\sim$0.85}\sym{b} & \cite{schoonbeek2024supcon} \\
& SupCon+ISIL (ResNet-34)   & MAP@R(+)       & \todo & \cite{schoonbeek2024supcon} \\
& Best contrastive (ResNet-34) & Error-Verif.\ AP & \texttt{$\sim$0.58}\sym{b} & \cite{lehman2024statediffnet} \\
& \popw{} (ours)            & F1@1           & \popwres & --- \\
& \popw{} (ours)            & Error-Verif.\ AP & \popwres & --- \\
\midrule
\multicolumn{5}{l}{\textit{Head pose (9-DoF)}} \\
& \textit{no published supervised baseline} & --- & --- & --- \\
& \popw{} (ours)            & Forward angular MAE ($^\circ$)  & \popwres & --- \\
& \popw{} (ours)            & Up angular MAE ($^\circ$)       & \popwres & --- \\
& \popw{} (ours)            & Position MAE (mm)               & \popwres & --- \\
\bottomrule
\end{tabular}

\vspace{0.4em}
\footnotesize
\sym{a} STORM-PSR uses a $\pm 3$-frame tolerance. We report \popw{} under both
$\pm 3$ and $\pm 5$ tolerances for fair comparison.
\sym{b} Estimated from the published figures (numbers not given in tabular form).
\end{table}

\FloatBarrier
% =====================================================================
\section{Per-task breakdown}
\label{sec:per-task}

The headline tables interleave tasks for compactness. The following
subsections split results by task, which is the comparison form most
reviewers reach for when checking a single capability.

\subsection{Activity recognition}
\label{sec:activity-breakdown}

\begin{table}[!htbp]
\centering
\small
\caption{Activity recognition: per-protocol breakdown across both datasets.
Empty cells are reserved for \popw{} results.}
\label{tab:activity-breakdown}
\renewcommand{\arraystretch}{1.15}
\begin{tabular}{l l c c c}
\toprule
\textbf{Dataset} & \textbf{Method} & \textbf{Top-1} & \textbf{Top-5} & \textbf{mcAP} \\
\midrule
\ikea{} (front, RGB)    & P3D~\cite{benshabat2021ikea}                 & 60.40 & \todo & \todo \\
\ikea{} (front, RGB+pose) & I3D~\cite{benshabat2021ikea}               & 64.15 & \todo & \todo \\
\ikea{} (all views)     & I3D~\cite{benshabat2021ikea}                 & 47.00 & \todo & \todo \\
\ikea{} (all views)     & PC3D~\cite{aganian2023ikea}                  & 80.20 & \todo & \todo \\
\ikea{} (online, csv)   & PTMA~\cite{xie2025ptma}                      & \todo & \todo & 84.47 \\
\ikea{} (online, cs)    & PTMA~\cite{xie2025ptma}                      & \todo & \todo & 86.99 \\
\ikea{} (online, cs)    & MiniROAD~\cite{xie2025ptma}                  & \todo & \todo & 80.84 \\
\ikea{}                 & \popw{} (ours)                               & \popwres & \popwres & \popwres \\
\midrule
\industreal{}           & MViTv2~\cite{schoonbeek2024industreal}       & 66.45 & 88.43 & \todo \\
\industreal{}           & \popw{} (ours)                               & \popwres & \popwres & \popwres \\
\bottomrule
\end{tabular}
\end{table}

\subsection{Object / state detection}
\label{sec:detection-breakdown}

\begin{table}[!htbp]
\centering
\small
\caption{Detection performance across both datasets. \industreal{} ASD is a 24-class
state recognition problem treated as detection by~\cite{schoonbeek2024industreal}.}
\label{tab:detection-breakdown}
\renewcommand{\arraystretch}{1.15}
\begin{tabular}{l l c c}
\toprule
\textbf{Dataset} & \textbf{Method} & \textbf{AP} & \textbf{AP@0.5} \\
\midrule
\ikea{}        & ResNeXt-101-FPN~\cite{benshabat2021ikea}   & 65.9 & 85.3 \\
\ikea{}        & Mask R-CNN~\cite{benshabat2021ikea}        & \todo & 78.9 \\
\ikea{}        & \popw{} (ours)                             & \popwres & \popwres \\
\midrule
\industreal{}  & YOLOv8m~\cite{schoonbeek2024industreal}    & \todo & 83.8 \\
\industreal{}  & \popw{} (ours)                             & \popwres & \popwres \\
\bottomrule
\end{tabular}
\end{table}

\subsection{Pose estimation}
\label{sec:pose-breakdown}

\begin{table}[!htbp]
\centering
\small
\caption{Pose results. The two datasets define pose differently: \ikea{} reports 2D body
keypoints (PCK), \industreal{} provides 9-DoF head pose with no published supervised baseline.
\popw{} establishes the first such baseline on \industreal{}.}
\label{tab:pose-breakdown}
\renewcommand{\arraystretch}{1.15}
\begin{tabular}{l l c}
\toprule
\textbf{Dataset / Pose} & \textbf{Method / Metric} & \textbf{Score} \\
\midrule
\multicolumn{3}{l}{\textit{\ikea{} --- 2D body pose, 17 keypoints}} \\
MaskRCNN-ft~\cite{benshabat2021ikea}    & PCK@10\,px         & 64.3 \\
MaskRCNN-ft~\cite{benshabat2021ikea}    & PCK@0.2            & 88.0 \\
\popw{} (ours)                          & PCK@10\,px         & \popwres \\
\popw{} (ours)                          & PCK@0.2            & \popwres \\
\midrule
\multicolumn{3}{l}{\textit{\industreal{} --- 9-DoF head pose}} \\
\popw{} (ours, baseline)                & Forward angular MAE ($^\circ$) & \popwres \\
\popw{} (ours, baseline)                & Up angular MAE ($^\circ$)      & \popwres \\
\popw{} (ours, baseline)                & Position MAE (mm)              & \popwres \\
\bottomrule
\end{tabular}
\end{table}

\subsection{Procedure Step Recognition (\industreal{})}
\label{sec:psr-breakdown}

\begin{table}[!htbp]
\centering
\small
\caption{PSR on \industreal{}. Tolerance window is the temporal slack allowed when matching
predicted step boundaries to ground truth. STORM-PSR is the strongest published baseline.}
\label{tab:psr-breakdown}
\renewcommand{\arraystretch}{1.15}
\begin{tabular}{l c c c c c}
\toprule
\textbf{Method} & \textbf{Tol.} & \textbf{Precision} & \textbf{Recall} & \textbf{F1} & \textbf{POS} \\
\midrule
B3 rule-based~\cite{schoonbeek2024industreal} & $\pm 5$ & 0.885 & 0.880 & 0.883 & \todo \\
STORM-PSR~\cite{schoonbeek2025storm}          & $\pm 3$ & \todo & \todo & 0.901 & 0.812 \\
\popw{} (ours)                                & $\pm 5$ & \popwres & \popwres & \popwres & \popwres \\
\popw{} (ours)                                & $\pm 3$ & \popwres & \popwres & \popwres & \popwres \\
\bottomrule
\end{tabular}
\end{table}

\FloatBarrier
% =====================================================================
\section{Efficiency comparison}
\label{sec:efficiency}

\popw{} is structurally heavier than purpose-built single-task models such as
PTMA (12.9\,M parameters) or MiniROAD (10.5\,M); we accept that trade.
The fair comparison is against a sequential pipeline of single-task baselines
that would be required to produce the same set of outputs --- an arrangement
in which throughput is bounded by the slowest stage.

\begin{table}[!htbp]
\centering
\small
\caption{Efficiency at native dataset resolution.
\popw{} reports two FPS numbers: \emph{batched} (single-frame forward, the standard reporting convention)
and \emph{streaming} (per-frame forward with cached temporal feature bank, the deployment-relevant number).
Where competitor papers do not report a number we use \todo. ``$^\dagger$'' denotes RTX 3060.}
\label{tab:efficiency}
\renewcommand{\arraystretch}{1.15}
\begin{tabular}{l l S[table-format=2.1] S[table-format=3.2] S[table-format=3.0] c}
\toprule
\textbf{Dataset} & \textbf{Method} &
{\textbf{Params (M)}} & {\textbf{GFLOPs}} &
{\textbf{FPS}} & \textbf{Hardware} \\
\midrule
\multicolumn{6}{l}{\textit{\ikea{} --- $640{\times}480$}} \\
\ikea{} & MiniROAD~\cite{xie2025ptma}              & 10.5 &   1.08 & 325 & --- \\
\ikea{} & PTMA~\cite{xie2025ptma}                  & 12.9 &   1.96 & 291 & --- \\
\ikea{} & ActionFormer (RGB)~\cite{gatedsrm2026}   & 27.7 &  83.28 &  21 & GTX\,1080 \\
\ikea{} & Gated SRM~\cite{gatedsrm2026}            & 33.5 & 121.09 &  16 & GTX\,1080 \\
\ikea{} & \popw{} (ours, batched)                  & {\popwres} & {\popwres} & {\popwres} & RTX\,3060$^\dagger$ \\
\ikea{} & \popw{} (ours, streaming)                & {\popwres} & {\popwres} & {\popwres} & RTX\,3060$^\dagger$ \\
\midrule
\multicolumn{6}{l}{\textit{\industreal{} --- $1280{\times}720$}} \\
\industreal{} & YOLOv8m~\cite{schoonbeek2024industreal}     & {\todo} & {\todo} & {\todo} & --- \\
\industreal{} & MViTv2~\cite{schoonbeek2024industreal}      & {\todo} & {\todo} & {\todo} & --- \\
\industreal{} & STORM-PSR~\cite{schoonbeek2025storm}        & {\todo} & {\todo} & {\todo} & --- \\
\industreal{} & \popw{} (ours, batched)                     & {\popwres} & {\popwres} & {\popwres} & RTX\,3060$^\dagger$ \\
\industreal{} & \popw{} (ours, streaming)                   & {\popwres} & {\popwres} & {\popwres} & RTX\,3060$^\dagger$ \\
\bottomrule
\end{tabular}
\end{table}

\begin{table}[!htbp]
\centering
\small
\caption{Verified benchmark papers with publicly available model weights/code for efficiency testing on user GPU.}
\label{tab:verified-models}
\renewcommand{\arraystretch}{1.15}
\begin{tabular}{l l l l}
\toprule
\textbf{Model} & \textbf{arXiv ID} & \textbf{Venue} & \textbf{Public Code} \\
\midrule
Ben-Shabat et al. (IKEA ASM) & \href{https://arxiv.org/abs/2007.00394}{2007.00394} & WACV 2021 & \href{https://github.com/ultralytics/ultralytics}{YOLOv8m Available} \\
Xie et al. (PTMA / MiniROAD) & \href{https://arxiv.org/abs/2508.17025}{2508.17025} & IEEE TMM 2025 & Code Not Released \\
Takami \& Fukuda (Gated SRM) & \href{https://www.mdpi.com/1424-8220/26/8/2454}{Sensors 2026} & MDPI 2026 & Code Not Released \\
Zhang et al. (ActionFormer) & \href{https://arxiv.org/abs/2202.07925}{2202.07925} & ECCV 2022 & \href{https://github.com/happyharrycn/actionformer\_release}{Available} \\
Schoonbeek et al. (IndustReal) & \href{https://arxiv.org/abs/2310.17323}{2310.17323} & WACV 2024 & \href{https://github.com/ultralytics/ultralytics}{YOLOv8m}, \href{https://github.com/facebookresearch/SlowFast}{MViTv2 via SlowFast} \\
Schoonbeek et al. (STORM-PSR) & \href{https://arxiv.org/abs/2510.12385}{2510.12385} & CVIU 2025 & \href{https://github.com/shaohsuanhung/STORM-PSR}{Available} \\
Schoonbeek et al. (SupCon+ISIL) & \href{https://arxiv.org/abs/2408.11700}{2408.11700} & IEEE RAL 2024 & Code Not Released \\
Lehman et al. (StateDiffNet) & \href{https://arxiv.org/abs/2408.12945}{2408.12945} & ECCV VISION 2024 & Code Not Released \\
\bottomrule
\end{tabular}
\end{table}

\vspace{0.5em}
\footnotesize
Note: Section~5 efficiency comparison uses YOLOv8m, MViTv2, STORM-PSR, ActionFormer, PTMA, MiniROAD, and Gated SRM. YOLOv8m code is from ultralytics (not the original authors). MViTv2 code is via PySlowFast~\cite{schoonbeek2024industreal}. ActionFormer code is at happyharrycn/actionformer\_release. PTMA and MiniROAD share the same arXiv paper (2508.17025) but have no public code. Gated SRM is a Sensors 2026 paper by Takami \& Fukuda with no public code.

\begin{table}[!htbp]
\centering
\small
\caption{Known benchmark metrics from verified papers for reference.}
\label{tab:verified-metrics}
\renewcommand{\arraystretch}{1.15}
\begin{tabular}{l l c c}
\toprule
\textbf{Model} & \textbf{Dataset} & \textbf{Metric} & \textbf{Value} \\
\midrule
SlowFast (Feichtenhofer et al.) & Kinetics-400 & Top-1 & 79.8\% \\
SlowFast (Feichtenhofer et al.) & AVA v2.2 & mAP & 28.3\% \\
TSM (Lin et al.) & Kinetics-400 & Top-1 & 76\% \\
Non-Local (Wang et al.) & Kinetics-400 & Top-1 & 77.7\% \\
ST-GCN (Yan et al.) & NTU RGB+D & Accuracy & 81.5\% \\
\bottomrule
\end{tabular}
\end{table}

\begin{table}[!htbp]
\centering
\small
\caption{Sequential single-task pipeline vs.\ \popw{}'s single forward pass on \industreal{}.}
\label{tab:multi-model}
\renewcommand{\arraystretch}{1.15}
\begin{tabular}{l c c c}
\toprule
\textbf{Pipeline} & \textbf{Params} & \textbf{GFLOPs} & \textbf{FPS} \\
\midrule
YOLOv8m + MViTv2 + STORM-PSR (sequential) & {\todo} & {\todo} & {\todo} \\
\popw{} (one forward pass)                & \popwres & \popwres & \popwres \\
\bottomrule
\end{tabular}
\end{table}

\FloatBarrier
% =====================================================================
\section{Target thresholds}

\begin{table}[!htbp]
\centering
\small
\caption{Sequential single-task pipeline vs.\ \popw{}'s single forward pass on \industreal{}.}
\label{tab:multi-model}
\renewcommand{\arraystretch}{1.15}
\begin{tabular}{l c c c}
\toprule
\textbf{Pipeline} & \textbf{Params} & \textbf{GFLOPs} & \textbf{FPS} \\
\midrule
YOLOv8m + MViTv2 + STORM-PSR (sequential) & {\todo} & {\todo} & {\todo} \\
\popw{} (one forward pass)                & \popwres & \popwres & \popwres \\
\bottomrule
\end{tabular}
\end{table}

\FloatBarrier
% =====================================================================
\section{Target thresholds}
\label{sec:targets}

For transparency we list the exact value \popw{} must beat on each metric to constitute
state-of-the-art. The thresholds are taken from the strongest published baseline.

\begin{table}[!htbp]
\centering
\small
\caption{\popw{} target thresholds for each headline metric. ``Beat'' means strictly greater
than (or strictly less than for error metrics) the listed value.}
\label{tab:targets}
\renewcommand{\arraystretch}{1.15}
\begin{tabular}{l l l c}
\toprule
\textbf{Dataset} & \textbf{Metric} & \textbf{Strongest baseline} & \textbf{Threshold} \\
\midrule
\ikea{}         & Object det.\ AP@0.5             & ResNeXt-101-FPN  & $>$\,85.30 \\
\ikea{}         & Object det.\ AP (COCO)          & ResNeXt-101-FPN  & $>$\,65.90 \\
\ikea{}         & Pose PCK@0.2                    & MaskRCNN-ft      & $>$\,88.00 \\
\ikea{}         & Pose PCK@10\,px                 & MaskRCNN-ft      & $>$\,64.30 \\
\ikea{}         & Activity Top-1 (front, RGB+pose) & I3D RGB+pose     & $>$\,64.15 \\
\ikea{}         & Activity Top-1 (all views)      & PC3D             & $>$\,80.20 \\
\ikea{}         & Activity mcAP (csv)             & PTMA             & $>$\,84.47 \\
\ikea{}         & Activity mcAP (cs)              & PTMA             & $>$\,86.99 \\
\ikea{}         & Temporal loc.\ \map{}@0.5       & Gated SRM        & $>$\,21.77 \\
\ikea{}         & Phase Acc@1.0                   & STEPs            & $>$\,37.02 \\
\ikea{}         & Temporal order Kendall's $\tau$ & STEPs            & $>$\,0.91 \\
\midrule
\industreal{}   & ASD \map{}@0.5                  & YOLOv8m          & $>$\,83.80 \\
\industreal{}   & Activity Top-1                  & MViTv2           & $>$\,66.45 \\
\industreal{}   & Activity Top-5                  & MViTv2           & $>$\,88.43 \\
\industreal{}   & PSR F1 ($\pm 3$ frames)         & STORM-PSR        & $>$\,0.901 \\
\industreal{}   & PSR POS ($\pm 3$ frames)        & STORM-PSR        & $>$\,0.812 \\
\industreal{}   & Assembly state F1@1             & SupCon+ISIL      & $>$ baseline \\
\industreal{}   & Error-verification AP           & GCA              & $>$ baseline \\
\industreal{}   & Head pose (9-DoF)               & none             & establish \\
\bottomrule
\end{tabular}
\end{table}

\FloatBarrier
% =====================================================================
\section{Reproducibility}
\label{sec:repro}

All competitor numbers are taken from the cited papers. Where a paper reports multiple
protocols (e.g.\ cross-subject, cross-view, cross-subject-view for online action detection
on \ikea{}~\cite{xie2025ptma}), we report the strongest. \popw{} is evaluated under each
listed protocol exactly; in particular,
PSR is evaluated at both $\pm 3$ and $\pm 5$ frame tolerances to permit direct comparison
against both STORM-PSR~\cite{schoonbeek2025storm} (which uses $\pm 3$) and the
\industreal{} paper's B3 baseline~\cite{schoonbeek2024industreal} (which uses $\pm 5$).
Final \popw{} numbers will be reported as mean$\pm$std across three random seeds
(42, 123, 7).

% =====================================================================
% Bibliography
% =====================================================================
\begin{thebibliography}{9}

\bibitem{benshabat2021ikea}
Y.~Ben-Shabat, X.~Yu, F.~Saleh, D.~Campbell, C.~Rodriguez-Opazo, H.~Li, and S.~Gould,
``The IKEA ASM Dataset: Understanding people assembling furniture through actions, objects and pose,''
in \emph{Proc.\ IEEE/CVF Winter Conf.\ on Applications of Computer Vision (WACV)}, 2021.
arXiv:2007.00394.

\bibitem{shah2023steps}
A.~Shah, B.~Lundell, H.~Sawhney, and R.~Chellappa,
``STEPs: Self-Supervised Key Step Extraction and Localization from Unlabeled Procedural Videos,''
in \emph{Proc.\ IEEE/CVF Int.\ Conf.\ on Computer Vision (ICCV)}, 2023.
arXiv:2301.00794.

\bibitem{aganian2023ikea}
D.~Aganian, M.~K\"ohler, S.~Baake, M.~Eisenbach, and H.-M.~Gross,
``How Object Information Improves Skeleton-based Human Action Recognition in Assembly Tasks,''
in \emph{Proc.\ Int.\ Joint Conf.\ on Neural Networks (IJCNN)}, 2023.
arXiv:2306.05844.

\bibitem{xie2025ptma}
L.~Xie, Y.~Tan, S.~Jing, H.~Lu, and K.~Zhang,
``Probabilistic Temporal Masked Attention for Cross-view Online Action Detection,''
\emph{IEEE Trans.\ on Multimedia}, 2025 (in press).
arXiv:2508.17025.

\bibitem{gatedsrm2026}
M.~Takami and T.~Fukuda,
``Confidence-Aware Gated Multimodal Fusion for Robust Temporal Action Localization in Occluded Environments,''
\textit{Sensors}, vol.~26, no.~8, article 2454, 2026.
doi: 10.3390/s26082454.

\textit{Note:} ActionFormer baseline results reported in this paper are from Zhang et al.~(ECCV 2022)~\cite{zhang2022actionformer}.

\bibitem{zhang2022actionformer}
C.-L.~Zhang, J.-X.~Wu, and Y.~Li,
``ActionFormer: Localizing Moments of Actions with Transformers,''
in \emph{Proc.\ European Conf.\ on Computer Vision (ECCV)}, 2022.
arXiv:2202.07925.
Code: \url{https://github.com/happyharrycn/actionformer\_release}.

\bibitem{schoonbeek2024industreal}
T.~J.~Schoonbeek, T.~Houben, H.~Onvlee, P.~H.~N.~de~With, and F.~van~der~Sommen,
``IndustReal: A Dataset for Procedure Step Recognition Handling Execution Errors
in Egocentric Videos in an Industrial-Like Setting,''
in \emph{Proc.\ IEEE/CVF Winter Conf.\ on Applications of Computer Vision (WACV)}, 2024.
arXiv:2310.17323.

\bibitem{schoonbeek2025storm}
T.~J.~Schoonbeek, S.-H.~Hung, J.~Kustra, P.~H.~N.~de~With, and F.~van~der~Sommen,
``STORM-PSR: Spatio-Temporal Reasoning for Procedure Step Recognition,''
\emph{Computer Vision and Image Understanding (CVIU)}, 2025.
arXiv:2510.12385.

\bibitem{schoonbeek2024supcon}
T.~J.~Schoonbeek, G.~Balachandran, H.~Onvlee, T.~Houben, S.-H.~Hung, J.~Kustra,
P.~H.~N.~de~With, and F.~van~der~Sommen,
``Supervised Representation Learning towards Generalizable Assembly State Recognition,''
\emph{IEEE Robotics and Automation Letters (RAL)}, 2024.
arXiv:2408.11700.

\bibitem{lehman2024statediffnet}
D.~Lehman, T.~J.~Schoonbeek, S.-H.~Hung, J.~Kustra, P.~H.~N.~de~With, and F.~van~der~Sommen,
``Find the Assembly Mistakes: Error Segmentation for Industrial Applications,''
in \emph{Proc.\ ECCV Vision-based InduStrial InspectiON (VISION) Workshop}, 2024.
arXiv:2408.12945.

\end{thebibliography}

\end{document}
